\begingroup
\setlength{\emergencystretch}{2em}
\capituloTesis{Metodología}

\seccionTesis{Material}

La metodología se organizó alrededor de registros georreferenciados de drenaje y pavimento correspondientes al corredor EMP.LI105 Pedregal–Censo–Cumbray, en Simbal, con el propósito de conservar una identificación común para las mediciones de continuidad longitudinal, abertura transversal, condición superficial y rugosidad. Los materiales de registro comprendieron fichas de inspección, formularios PCI, formatos de perfiles longitudinales, inventario de estructuras y fotografías con referencia espacial; a través de esos medios se definieron las observaciones necesarias para construir cuatro indicadores cuantitativos por unidad estructura–sector. La distinción entre fuentes, instrumentos y soportes de procesamiento permitió relacionar cada dato con su origen y con la regla operacional correspondiente, conforme al criterio de sistematización de las observaciones que describen Hernández Sampieri et al. (2014) y Bernal Torres (2016).

La información de partida incluyó la identificación del proyecto vial, los cuatro tramos del corredor, las progresivas de los pases de agua, las referencias geométricas de las secciones y las características constructivas disponibles. Esos antecedentes se emplearon para delimitar unidades y establecer la organización del recorrido; la medición de los indicadores, por su parte, se definió mediante formatos diferenciados para los márgenes, los extremos de las estructuras, las superficies de inspección PCI y las huellas longitudinales. Se preservó en cada registro el código de unidad, tramo, progresiva, fecha y tipo de observación, de manera que una anotación de drenaje pudiera identificarse junto con el sector de pavimento al que correspondía.

Los equipos previstos en el protocolo comprendieron wincha y regla para longitudes y dimensiones de aberturas, un dispositivo de posicionamiento para asociar coordenadas, elementos de referencia métrica para fotografías y un sistema estático de contacto para obtener perfiles de elevación. El dispositivo de perfilado se distinguió de los programas destinados al procesamiento porque su función correspondía a la adquisición de datos, mientras que las operaciones de cálculo se desarrollaron a partir de los registros obtenidos y de los algoritmos establecidos para los indicadores funcionales; ASTM International (2018) proporcionó el procedimiento de inspección del PCI y ASTM International (2021), el cálculo del IRI a partir del perfil longitudinal.

Para organizar la información se utilizó una estructura de hojas con registros de drenaje longitudinal por margen, drenaje transversal por extremo, unidades de inspección PCI y recorridos de rugosidad por huella, además de una tabla de integración con una fila por sector. Las hojas conservaron los numeradores y denominadores de los porcentajes, las áreas de las superficies inspeccionadas y los recorridos asociados con el perfil, permitiendo comprobar los valores derivados antes del análisis estadístico. La base disponible para el ejercicio de redacción y sustentación corresponde a un conjunto de escenarios de práctica, por lo que su consistencia interna no constituye por sí misma acreditación de observaciones de campo efectivamente ejecutadas.

Como recursos bibliográficos se utilizaron las prácticas ASTM correspondientes y el manual de drenaje del MTC, mientras que la revisión metodológica consideró las especificaciones de confiabilidad de McGraw y Wong (1996), la aleatorización espacial de Wagner y Dray (2015) y el procedimiento de ajuste de pruebas múltiples de Holm (1979). Estos documentos aportaron criterios de medición y análisis, sin reemplazar los registros que debían originarse en las unidades de estudio.

\seccionTesis{Población}

La población analítica del estudio se delimitó mediante veinte unidades estructura–sector asociadas con los veinte pases de agua identificados en el corredor EMP.LI105 Pedregal–Censo–Cumbray, distrito de Simbal, provincia de Trujillo, departamento de La Libertad, durante el periodo de evaluación de 2026. El corredor comprende una longitud documentada de 3\,192,16 m distribuida en cuatro tramos, cuyos sistemas de progresivas se mantuvieron separados para evitar que una misma referencia numérica identificara ubicaciones distintas. La unidad de análisis quedó constituida por el pase transversal y la superficie vial que se le asignó, mientras que las unidades de observación incluyeron márgenes de conducción, extremos de alcantarilla, subunidades de evaluación superficial y huellas de perfilado.

La delimitación espacial tomó como punto inicial la progresiva de cada pase y extendió el sector hasta cincuenta metros hacia ambos lados, cuando existía longitud disponible; en los extremos de tramo, el intervalo se recortó en su límite físico y, cuando dos sectores se aproximaban hasta superponerse, el punto medio entre las progresivas de las estructuras determinó la separación. Esa regla generó intervalos exclusivos, evitando asignar una misma superficie a dos unidades analíticas y conservando una estructura de datos compatible con la comparación entre sectores. Las mediciones de los dos márgenes se asociaron al mismo sector, al igual que las áreas interiores de entrada y salida del paso de agua.

La población comprendió sectores con antecedentes constructivos diferenciados, incluida la referencia de recapeo en los primeros 325 m del tramo 01, condición que se registró como atributo descriptivo para interpretar la variabilidad superficial y longitudinal. No se consideraron como unidades independientes las múltiples fotografías, los recorridos de repetición del perfilado ni las subunidades PCI: cada uno representó un componente o registro subordinado a la misma unidad estructura–sector, por lo que su acumulación no aumentó el número de observaciones disponibles para contrastar las hipótesis.

La delimitación temporal se restringió a 2026 y correspondió a una campaña de observación transversal; las diferencias entre fechas de registro se controlaron mediante el código y la ventana de medición de cada unidad. La población así definida constituyó un conjunto finito de referencia para la investigación, de modo que las conclusiones se circunscribieron al entorno de los pases identificados y no representaron automáticamente toda la superficie del corredor ni otras carreteras de características semejantes.

\seccionTesis{Muestra}

El estudio estableció un procedimiento censal sobre las veinte unidades que integraban la población analítica, con tamaño previsto y estructura de base \(n=N=20\). La decisión se fundamentó en la disponibilidad de un conjunto finito y completamente identificado de pases de agua, cada uno asociado con un sector de pavimento y con mediciones que debían compartir límites espaciales; por consiguiente, no se aplicó una fórmula de tamaño muestral ni se efectuó selección aleatoria de unidades dentro del corredor. Hernández Sampieri et al. (2014) distinguen el marco de unidades de la estrategia de selección, distinción relevante porque el carácter censal de la cobertura no convierte en probabilística la elección previa del corredor.

El criterio de inclusión comprendió los pases pertenecientes a los cuatro tramos definidos, con progresiva identificable y un sector delimitable mediante la regla de cincuenta metros y sus respectivos recortes. Para completar el registro, la unidad debía contar con información evaluable de longitud longitudinal, aberturas transversales y cobertura de las mediciones funcionales; cuando un extremo de alcantarilla no podía observarse, esa condición debía registrarse explícitamente y el indicador correspondiente no debía sustituirse por un valor inferido. Los criterios de exclusión se referían a elementos ajenos al inventario de pases y a observaciones fuera de la delimitación del corredor o del periodo establecido.

La diferencia entre ausencia de una estructura, abertura físicamente ocupada y dato faltante se preservó en la codificación, puesto que una obstrucción completa puede recibir valor de área libre cero cuando se verificó su condición, mientras que la imposibilidad de medir no constituye evidencia de ausencia de abertura. Para los denominadores geométricos nulos o no verificables se definió un resultado no calculable, sujeto a revisión del registro de origen; la misma regla de trazabilidad se extendió a sectores sin perfiles suficientes o con una superficie PCI no conciliada. El inventario censal se mantuvo como referencia, aunque el número de pares válidos para cada contraste debía informarse según la disponibilidad efectiva de ambos indicadores.

El procedimiento de selección de los registros operó, por tanto, en dos niveles: identificación exhaustiva de las veinte unidades pertenecientes al caso y control de integridad de los datos empleados en cada análisis. La cantidad de fotografías, subunidades superficiales y repeticiones de perfil no se utilizó como multiplicador de \(n\), pues hacerlo habría introducido seudorreplicación en las estimaciones de asociación. El tamaño reducido del conjunto y la vecindad geográfica entre estructuras constituyeron restricciones relevantes para la precisión de los coeficientes y justificaron examinar explícitamente la dependencia espacial durante la contrastación.

